\section{Classification Results}

In this section, I present the results of the filtering and classification procedures applied to the aggregated database. The analysis is structured on a per-repository basis in order to provide a clear and systematic overview of the observed patterns. Each repository-specific subsection contains a set of standardized outputs designed to facilitate comparison across repositories and to support interpretation of the classification results.

More specifically, each subsection includes the following components:

\begin{enumerate}
    \item A histogram and a pie chart illustrating the distribution of project types identified within the repository. These visualizations provide an overview of the relative frequency and proportional composition of the different project categories.

    \item A histogram showing the distribution of primary classes detected in the repository. This allows for a detailed examination of the most prevalent classification outcomes and highlights dominant patterns in the data.

    \item A rank-ordered table listing the identified classes according to their frequency. This table provides a compact summary of class prevalence and supports direct comparison between classes in terms of their occurrence.
\end{enumerate}

\subsection{Zenodo}

\subsection{Harvard Dataverse}

\subsection{DANS}

\subsection{Dryad}

\subsection{Open Data Uni Halle}

\subsection{Finnish Social Science Data Archive}

\subsection{UK Data Service}

\subsection{Syracuse Qualitative Data Repository}

\subsection{SIKT}

\subsection{CESSDA}

\subsection{ICPSR}


\subsection{ADA}

\subsection{Harvard Murray Archive}

\subsection{DataverseNO}

\subsection{Columbia Oral History Archive}

\subsection{HIHSN}

\subsection{SADA}

\subsection{N/A}

\subsection{AUSSDA}
